\chapter{Мост Эйнштейна-Розена}\label{ch:2}
% \cite{einsteinParticleProblemGeneral1935, ellisEtherFlowDrainhole1973, dokuchaevGENERALIZEDEINSTEINROSEN2024}

\section{Классическая формулировка}\label{sec:2.1}
Фундаментальная концепция пространственно-временного моста, ныне известного в научной литературе как мост Эйнштейна-Розена (ЭР-мост), была впервые сформулирована в 1935 году в результате совместных исследований Альберта Эйнштейна и Натана Розена \cite{einsteinParticleProblemGeneral1935}. Данная работа является результатом попыток геометризации физики, направленных на построение единой теории поля, где элементарные частицы рассматривались бы как топологические особенности пространства-времени.

В своей статье авторы проводят тщательный анализ решения Шварцшильда для сферически-симметричного гравитационного поля в рамках ОТО. Рассматривая решение Шварцшильда \cref{eq:Schw_metric} они обнаружили, что сингулярность при \(r = 2M\) может быть интерпретирована как топологическая <<горловина>> , соединяющая два асимптотически плоских пространства-времени через минимальную поверхность \(S^2\) \cite{ellisEtherFlowDrainhole1973}.
Следует подчеркнуть, что в первоначальной интерпретации Эйнштейна и Розена данная конструкция рассматривалась исключительно как математическая модель элементарной частицы (в частности, электрона), где:
\begin{enumerate}[beginpenalty=10000]
	\item горловина моста моделировала протяженную структуру частицы;
	\item два соединенных пространства-времени представляли <<внутреннюю>> и <<внешнюю>> геометрию;
	\item топологический переход заменял точечную сингулярность.
\end{enumerate}
Современный термин <<мост Эйнштейна-Розена>> был введен значительно позже.

Теоретическая значимость работы Эйнштейна и Розена заключается в фундаментальном переосмыслении природы сингулярностей в решениях уравнений ОТО, продемонстрировавшем возможность их интерпретации как топологических особенностей пространства-времени, а не математических патологий. Введя оригинальную систему координат, авторы впервые показали, что решение Шварцшильда может описывать не точечную сингулярность, а топологический переход между двумя асимптотически плоскими пространствами, тем самым заложив концептуальные основы для современных исследований кротовых нор и экзотических топологий в релятивистской гравитации. Хотя первоначальная модель оказалась физически нереализуемой в рамках классической ОТО из-за динамической неустойчивости и необходимости экзотической материи, она стимулировала развитие теоретической физики, включая проблему регуляризации сингулярностей, исследования глобальной структуры пространства-времени и квантовые подходы к гравитации, сохранив методологическую ценность как первый пример нетривиальной топологической интерпретации решений уравнений Эйнштейна.

Концепция моста Эйнштейна-Розена получила своё развитие в современной теории кротовых нор — гипотетических топологических тоннелей в пространстве-времени. В отличие от исходной модели, которая представляла собой неустойчивое соединение двух вселенных через шварцшильдовскую сингулярность, современные кротовые норы рассматриваются как потенциально проходимые при определённых условиях структуры. Ключевое отличие заключается в том, что если мост Эйнштейна-Розена был неотъемлемой частью решения для чёрной дыры и требовал наличия сингулярности, то кротовые норы в современной трактовке \cite{morris1988wormholes}  представляют собой самостоятельные топологические объекты, которые могут существовать и без связи с чёрными дырами. Для их стабилизации требуется экзотическая материя, нарушающая энергетические условия, что стало предметом активных исследований в квантовой теории поля и теориях квантовой гравитации.

Современные исследования кротовых нор как развития идеи Эйнштейна-Розена охватывают несколько перспективных направлений: в рамках теории струн предложены модели стабильных кротовых нор, поддерживаемых космическими струнами или восьмимерными бранами \cite{maldacena2021humanly}; Квантовая гравитация предлагает механизмы стабилизации за счёт эффектов квантовой пены пространства-времени \cite{wheeler1957nature, misner1957classical} или петлевых квантовых поправок \cite{ashtekar2018quantum}; особый интерес представляет голографическая интерпретация, где кротовые норы рассматриваются как пространственное проявление квантовой запутанности (ER=EPR, \cite{maldacenaCoolHorizonsEntangled2013}). Эти исследования  развивают первоначальную идею Эйнштейна-Розена, а так же связывают её с другими фундаментальными концепциями современной физики, открывая новые перспективы для понимания природы пространства-времени и квантовой гравитации.

\section{Обобщенный Мост Эйнштейна-Розена}\label{sec:2.2}

Мост ЭР может возникать в других решениях уравнений Эйнштейна, таких как метрика Керра-Ньюмена, несмотря на то, что является физически непроходимым в геометрии Шварцшильда. На основании этого было введено понятие обобщенного моста Эйнштейна-Розена внутри черной дыры, определенное как времени-подобное соединение между двумя вселенными с областями асимптотически минковских пространственно-временных бесконечностей.

Первоначальная идея о серии подобных мостов между вселенными в метрике Керра принадлежит Бойеру и Линдквисту \cite{boyer1967maximal}. Односторонний мост в метриках Рейсснера–Нордстрёма и Керра обсуждается в \cite{carrollSpacetimeGeometryIntroduction2014}. Также есть указания на физические препятствия для существования такого моста между вселенными, которые могут быть связаны с различными типами нестабильностей (включая квантовые), обсуждаемых, например, в \cite{gursel1979final, simpsonInternalInstabilityReissnerNordstrm1973}. Однако проблема остается открытой: симметричное геодезическое движение, связанные и несвязанные орбиты, а также возможность прохождения физического наблюдателя через мост Рейсснера–Нордстрёма и Керра были недавно вновь проанализированы в \cite{abramson2020momentum} и \cite{dyson2023kerr} соответственно.

\subsection{Основы метрики Керра–Ньюмена}\label{sub:2.2.1}

Рассмотрим Мост ЭР в наиболее общей из представленных в разделе \cref{sub:1.1.1} метрике — метрике Керра-Ньюмена, являющейся точным решением уравнений Эйнштейна для вращающейся и электрически заряженной черной дыры, которая имеет вид:
\begin{equation}
	\label{eq:KN}
	ds^2=-\frac{\Delta}{\Sigma}[dt-a\sin^2\theta d\varphi]^2+\frac{\sin^2\theta}{\Sigma}[(r^2+a^2)d\varphi -adt]^2  +\frac{\Sigma}{\Delta}dr^2+\Sigma d\theta^2\!,
\end{equation}
где \((r,\theta,\phi)\) — сферические координаты, а \(t\) — время статического удаленного наблюдателя на радиальной бесконечности. В этой метрике
\begin{subequations}
    \label{eq:DS}
    \begin{gather}
        \Delta = r^2-2Mr+a^2+q^2 \label{eq:Delta} \\
        \Sigma= r^2+a^2\cos^2\theta \label{eq:Sigma}
    \end{gather}
\end{subequations}
\(M\) — масса черной дыры, \(q\) — ее электрический заряд, \(a= J/M\) — удельный угловой момент (спин) черной дыры. Два корня уравнения \(\Delta=0\) дают радиусы \(r_±\) для положения горизонта событий и внутреннего горизонта Коши черной дыры:
\begin{subequations}
    \label{eq:rads}
    \begin{gather}
        r_+=M+\sqrt{M^2-a^2-q^2}, \label{r+} \\
        r_-=M-\sqrt{M^2-a^2-q^2}, \label{r-}
    \end{gather}
\end{subequations}
В данной работе мы рассматриваем метрику Керра с горизонтом событий, соответствующим условию \(a^2+q^2<M^2\) (случай наличия обоих горизонтов).

Для упрощения уравнений и представления графиков в дальнейшем будет использована система единиц, в которой \(M=1\), и соответствующие безразмерные параметры: радиус \(r\Rightarrow r/M\), время \(t\Rightarrow t/M\), спин черной дыры \(a\Rightarrow a/M\) и заряд \(q\Rightarrow q/M\).

В метрике Керра–Ньюмена существуют следующие интегралы движения для пробных частиц \cite{carterGlobalStructureKerr1968}: \(\mu\) — масса частицы, \(E\) — ее полная энергия, \(L\) — азимутальный угловой момент и \(Q\) — постоянная Картера, связанная с неэкваториальным движением частицы.

Соответствующие уравнения движения пробной частицы в метрике Керра–Ньюмена в дифференциальной форме имеют вид \cite{carterGlobalStructureKerr1968,thorne2000gravitation,chandrasekhar1998mathematical}:
\begin{subequations}
    \label{eq:EoM}
    \begin{gather}
	\Sigma \frac{dr}{d\tau}=\sqrt{R}, \label{dr} \\
	\Sigma \frac{d\theta}{d\tau} = \sqrt{\Theta}, \label{dtheta}\\
	\Sigma \frac{d\phi}{d\tau}=-(aE-\frac{L}{\sin^2\theta}) +\frac{a}{\Delta}P, \label{dvarphi}\\
	\Sigma \frac{dt}{d\tau} = -a(aE\sin^2\theta-L)+(r^2+a^2)\frac{P}{\Delta}. \label{dt}\\
    \end{gather}
\end{subequations}
Здесь
\begin{equation}
	P=E(r^2+a^2)-a L+\epsilon qr, \label{P}
\end{equation}
\(\tau\) — собственное время массивной пробной частицы или аффинный параметр вдоль траектории безмассовой частицы (\(\mu=0\)), такой как фотон, \(\epsilon\) — электрический заряд пробной частицы. Эффективный радиальный потенциал \(R(r)\) задается выражением:
\begin{equation}
	R(r) = P^2-\Delta[\mu^2r^2+(L-aE)^2+Q], \label{Rr}
\end{equation}
а эффективный полярный потенциал \(\Theta(\theta)\) имеет вид:
\begin{equation}
	\Theta(\theta) = Q-\cos^2\theta[a^2(\mu^2-E^2)+L^2\sin^{-2}\theta]. \label{Vtheta}
\end{equation}

Траектории массивных частиц \((\mu\neq0)\) зависят от трех параметров: \(\gamma=E/\mu\), \(\lambda=L/\mu\) и \(Q/\mu^2\). В то же время траектории безмассовых частиц, таких как фотоны (нулевые геодезические), зависят только от двух параметров: \(\lambda\) и \(Q\).

Нетривиальной особенностью вращающейся черной дыры Керра \((a\neq0)\) является существование так называемой эргосферы \cite{thorne2000gravitation,chandrasekhar1998mathematical,waldGeneralRelativity1984} с внешней границей
\begin{equation}\label{ES}
	r_{ES}(\theta) =1+\sqrt{1-q^2-a^2\cos^2\theta}.
\end{equation}
Внутри эргосферы любая пробная частица увлекается во вращение вокруг черной дыры с бесконечным азимутальным <<наматыванием>> при приближении к горизонту событий (отметим, что эффект <<наматывания>> также обсуждался в \cite{dyson2023kerr, kling2019examining}).

%В следующих разделах будет описана симметрия и асимметрия движения пробных частиц в гравитационном поле черной дыры Керра–Ньюмена. Мы используем уравнения движения \cref{dr,dtheta,dvarphi,dt} в наших аналитических и численных расчетах геодезических траекторий пробных частиц \cite{babichev2013black, dokuchaev2015physical, dokuchaev2017gravitational, dokuchaev2019event_a, dokuchaev2019event_b, dokuchaevSeeInvisibleImage2019}.

\subsection{Односторонний мост Эйнштейна–Розена внутри черной дыры}\label{sub:2.2.2}

%\subsection{Траектории движения пробных частиц}\label{sub:2.2.2}

Траектории пробных частиц в метрике Керра-Ньюмена могут быть разделены на 3 класса:
\begin{enumerate}
	\item над горизонтом событий,
	\item пересекающие горизонт событий и горизонт Коши,
	\item под горизонтом Коши.
\end{enumerate}
Устойчивые траектории в области между внешним и внутренним горизонтами не существуют, так как в этой области движение всегда направлено в сторону только уменьшения (или только увеличения) радиуса.

Для анализа траекторий движения частиц под горизонтом событий (прохождение Моста ЭР) была выбрана метрика Рейснера-Нордстрема (частный случай метрики Керра-Ньбмена со спином \(a=0\)), однако аналогичные выводы могут быть обобщены и на случай ненулевого спина.

Электрический заряд черной дыры составляет \(q=0.99\). Пробная частица с электрическим зарядом \(\epsilon=-1.5\) периодически обращается вокруг черной дыры с орбитальными параметрами \(\gamma=0.5\), \(\lambda=0.5\), соответствующими максимальному радиусу (апогею) \(r_{ max}=1.65\) и минимальному радиусу (перигей) \(r_{ min}=0.29\) в размерных единицах.

\begin{figure}[ht]
	\centerfloat{
	\hfill
	\subcaptionbox[List-of-Figures entry]{Траектории, пересекающие горизонт\label{fig:2D_traj_a}}{
		\includegraphics[angle=0,width=0.45\linewidth]{2D_traj_a}}
	\hfill
	\subcaptionbox{Траектории под горизонтом\label{fig:2D_traj_b}}{
		\includegraphics[angle=0,width=0.45\linewidth]{2D_traj_b}}
	\hfill
	}
	\legend{Сплошные окружности — горизонт событий и горизонт Коши. Пунктирные окружности — перигей и апогей}
	\caption[Двумерные траектории движения пробных частиц]{Двумерные траектории движения пробных частиц в метрике Рейсснера-Нордстрёма}\label{fig:2D_traj}
\end{figure}

Численное решение уравнений \cref{eq:EoM} с использованием встроенных методов в программный пакет Mathematica позволяет получить траектории движения пробных частиц в зависимости от их собственного времени. Примеры полученных траекторий в экваториальной плоскости представлены на рис. \cref{fig:2D_traj}: на рис. \cref{fig:2D_traj_a} изображена траектория движения в поле черной дыры с зарядом , пересекающая горизонт событий и горизонт Коши, в то время как на рис. \cref{fig:2D_traj_b} изображена траектория, полностью находящаяся под горизонтом Коши.

На рис. \cref{fig:2D_traj_a} электрический заряд черной дыры составляет \(e=0.99\), а пробная частица с зарядом \(\varepsilon=-1.5\) периодически обращается вокруг черной дыры с параметрами \(\gamma=1.7\), \(\lambda=0.1\), соответствующими максимальному радиусу \(r_{ max}=0.75\) и минимальному радиусу \(r_{ man}=0.09\). 

На рис. \cref{fig:2D_traj_b} электрический заряд черной дыры составляет \(q=0.99\). Пробная частица с электрическим зарядом \(\epsilon=-1.5\) периодически обращается вокруг черной дыры с орбитальными параметрами \(\gamma=0.5\), \(\lambda=0.5\), соответствующими максимальному радиусу (апогею) \(r_{ max}=1.65\) и минимальному радиусу (перигей) \(r_{ min}=0.29\).

Хотя на первый взгляд кажется, что картина движения геометрически абсолютно симметрична, так как путешествие начинается в апогее \(r_a\) из точки \(A\), достигает перигея \(r_p\) и возвращается в апогей \(r_a\) в точке \(B\) за конечное собственное время, демонстрируя абсолютную геометрическую симметрию, отметим ключевой момент, не отображенный на диаграмме: апогей \(r_a\) в точке \(B\) находится не в исходной пространственной области (вселенной), а в другой, отделенной от первоначальной, что ясно видно на диаграмме Картера-Пенроуза  на рис.~\cref{fig:Pen_diag}. 

\begin{figure}[ht]
	\centerfloat{
	\hfill
	\includegraphics[angle=0,width=0.5\linewidth]{Pen_diag} 
	\hfill
	}
	\caption[Диаграмма Картера-Пенроуза]{Диаграмма Картера-Пенроуза для сферически симметричной черной дыры Рейсснера-Нордстрёма с электрическим зарядом \(q=0.99\)} \label{fig:Pen_diag}
\end{figure}

Диаграмма, изображенная на рис.~\cref{fig:Pen_diag} позволяет наглядно просмотреть причинную структуру всего пространства-времени. Малиновой кривой изображена траектория наблюдателя, двигающегося из точки \(A\) в \(R_+\)-области в сторону центра черной дыры. После пересечения горизонта событий и горизонта Коши, оказавшись вблизи сингулярности черной дыры при \(r=0\) наблюдатель ускоряется в противоположную от сингулярности сторону и избегает приливного разрушения на малых радиусах. В результате он достигает точки \(B\) другой бесконечной асимптотически плоской области пространства.
С чисто геометрической точки зрения эта диаграмма симметрична как слева-направо, так и сверху-вниз. Напротив, с физической или пространственно-временной точки зрения эта диаграмма абсолютно асимметрична из-за неумолимого течения времени вверх не только на данной диаграмме, но и во всей Вселенной. Точнее, это означает, что в ОТО все объекты могут двигаться только внутри световых конусов, направленных вверх (под углом \(\pm45^\circ\) относительно направления вверх).

Такое явление было выше определено как Обобщенный Мост ЭР. В данном случае он, очевидно, является проходимым, но только в одну сторону. Следовательно в рассматриваемой метрике имеется проходимый односторонний Мост ЭР.

В тоже время траектория под горизонтом Коши, как на рис. \cref{fig:2D_traj_b}, является полностью симметричной как геометрически, так и физически: траектория представляет собой периодическое орбитальное движение вокруг центральной сингулярности черной дыры Рейсснера-Нордстрёма внутри горизонта Коши \(r_-\). Асимметричный Мост ЭР необходим только для проникновения в экзотическую область при \(0<r<r_-\), где существуют почти стабильные периодические орбиты для пробных частиц, очень похожие на периодические орбиты вне горизонта событий \(r_+\). Важно отметить, что периодическое движение пробной частицы ограничено во времени из-за потерь энергии на неизбежное излучение гравитационных волн.

\begin{figure}[ht]
	\centerfloat{
	\hfill
	\includegraphics[angle=0,width=0.6\linewidth]{insert_diag} 
	\hfill
	}
	\caption[Диаграмма вложения]{Диаграмма вложения для черной дыры Рейсснера-Нордстрёма}
	\label{fig:insert_diag}
\end{figure}

Мост ЭР может быть проиллюстрирован при помощи диаграммы вложения, представленной на рис. \cref{fig:insert_diag}. Этот мост соединяет две асимптотически плоские вселенные подобно туннелю кротовой норы  \cite{morris1988wormholes}, но только с односторонним движением от начальной точки \(A\) к конечной \(B\). Геометрическая симметрия этой диаграммы вложения также обманчива: фактически она демонстрирует асимметричную пространственно-временную природу одностороннего моста Эйнштейна-Розена (с учетом движения внутри светового конуса).

\subsection{Симметрия и асимметрия траекторий пробных частиц вблизи вращающейся черной дыры}\label{sub:2.2.3}

Аналогично были получены траектории движения пробных частиц в метрике вращающейся черной дыры Керра (\(a=0.9982\)), регистрируемые удаленным наблюдателем, расположенным на бесконечности. Эти траектории отличаются тем, что никогда не пересекают горизонт событий, а бесконечно долго <<наматываются>> на него.

\begin{figure}[ht]
	\centerfloat{
	\hfill
	\subcaptionbox[List-of-Figures entry]{\(\gamma=0.85\), \(\lambda=1.7\), \(Q=1\) \label{fig:traj_part_1}}{
		\includegraphics[angle=0,width=0.55\linewidth]{traj_part_1}}
	\hfill
	\subcaptionbox{\(\gamma=0.85\), \(\lambda=1.7\), \(Q=1\) \label{fig:traj_part_2}}{
		\includegraphics[angle=0,width=0.4\linewidth]{traj_part_2}}
	\hfill
	}
	\caption[Траектории пробной частицы]{Траектория пробной частицы (\(\mu\neq0\)) , падающей в быстро вращающуюся черную дыру Керра с спином \(a=0.9982\)}
	\label{fig:traj_part}
\end{figure}

\begin{figure}[ht]
\centerfloat{
	\hfill
	\subcaptionbox[List-of-Figures entry]{\(\lambda=2\), \(Q=1\) \label{fig:traj_ph_1}}{
		\includegraphics[angle=0,width=0.4\linewidth]{traj_ph_1}}
	\hfill
	\subcaptionbox{\(\lambda=-1.493\), \(Q=12.99\) \label{fig:traj_ph_2}}{
		\includegraphics[angle=0,width=0.45\linewidth]{traj_ph_2}}
	\hfill
	}
	\caption[Траектории фотонов]{Траектории фотонов с различными орбитальными параметрами, падающие в быстро вращающуюся черную дыру Керра со спином \(a=0.9982\)}
	\label{fig:traj_ph}
\end{figure}

Полученные траектории, представленные на рис. \cref{fig:traj_ph_1, fig:traj_ph_2, fig:traj_part_1,fig:traj_part_2}, демонстрируют как симметричные, так и асимметричные особенности траекторий массивных и безмассовых частиц, падающих во вращающуюся черную дыру Керра. Малиновые стрелки показывают направление вращения черной дыры согласно правилу буравчика. Разноцветные кривые на рис. \cref{fig:traj_ph_1, fig:traj_ph_2} представляют геодезические траектории безмассовых пробных частиц, таких как фотоны (\(\mu=0\)), рассчитанные численно с использованием уравнений движения \cref{eq:EoM}. При приближении к черной дыре траектории всех частиц, как массивных, так и безмассовых, бесконечно <<наматываются>> на горизонт событий в направлении вращения черной дыры на фиксированных широтах, что является проявлением симметрии поведения всех траекторий, падающих во вращающуюся черную дыру. В то же время направление вращения черной дыры является соответствующим проявлением асимметрии гравитационного поля метрики Керра.

Пробная частица на рис. \cref{fig:traj_part_1} начинает движение из верхнего полушария далеко от черной дыры, внутри эргосферы \cref{ES} частица движется по спирали вблизи горизонта событий черной дыры выше экваториальной плоскости. Пробная частица на рис. \cref{fig:traj_part_2} начинает движение из нижнего полушария далеко от черной дыры. Внутри эргосферы планета <<наматывается>> на горизонт событий черной дыры выше экваториальной плоскости. 


Фотон на рис. \cref{fig:traj_ph_1} падает в быстро вращающуюся черную дыру Керра с \(a=0.9982\) и <<наматывается>> на горизонт событий ниже экваториальной плоскости. 
Важно подчеркнуть, что на больших расстояниях от черной дыры пробная частица с отрицательным азимутальным угловым моментом (рис. \cref{fig:traj_ph_2}) вращается в направлении, противоположном вращению черной дыры, однако при приближении к черной дыре (внутри эргосферы) она вынуждена вращаться в том же направлении, что и черная дыра.

%Мы начинаем раскрывать универсальность увлекательных явлений симметрии и асимметрии вне и внутри черных дыр, используя диаграммы Картера–Пенроуза \cite{thorne2000gravitation,chandrasekhar1998mathematical,waldGeneralRelativity1984}, описывающие глобальную структуру пространства-времени черных дыр. Явным проявлением симметрии этой глобальной структуры являются бесконечные пространственные объемы как вне, так и внутри горизонта событий черной дыры. На рис. \cref{fig:Pen_diag} представлена соответствующая диаграмма Картера–Пенроуза для черной дыры Рейсснера–Нордстрёма — частного сферически { симметричного} случая черной дыры Керра–Ньюмена без вращения (\(a=0\), но \(q\neq0\)). С чисто геометрической точки зрения эта диаграмма симметрична как слева-направо, так и сверху-вниз. Напротив, с физической или пространственно-временной точки зрения эта диаграмма абсолютно асимметрична из-за неумолимого течения времени вверх не только на этой диаграмме, но и во всей Вселенной. Точнее, это означает, что в общей теории относительности все объекты могут двигаться только внутри световых конусов, направленных вверх (под углом \(\pm45^\circ\) относительно направления вверх). Направленный вверх световой конус представляет собой неумолимую { асимметрию} мира.
%
%Двумерное представление путешествия через внутренности черной дыры Рейсснера–Нордстрёма с использованием моста Эйнштейна–Розена показано на рис. \cref{fig:2D_traj_a}. Электрический заряд черной дыры составляет \(e=0.99\). Пробная планета (или космический корабль) с электрическим зарядом \(\varepsilon=-1.5\) периодически обращается вокруг черной дыры с орбитальными параметрами \(\gamma=0.5\), \(\lambda=0.5\), соответствующими максимальному радиусу (апогею) \(r_{ max}=1.65\) и минимальному радиусу (перигей) \(r_{ min}=0.29\) в размерных единицах.
%
%Периодические геодезические траектории планеты (малиновые кривые на рис. \cref{fig:2D_traj_a} и рис. \cref{fig:2D_traj_b}) были рассчитаны численно с использованием уравнений движения \cref{eq:EoM} для массивных пробных частиц (\(\mu\neq0\)). Отметим, что периодическое движение пробной планеты ограничено во времени из-за потерь энергии на неизбежное излучение гравитационных волн.
%
%Изображение на рис. \cref{fig:2D_traj_a} геометрически абсолютно симметрично или, другими словами, полностью или идеально { симметрично}. Геодезические траектории пробной планеты (\(\mu\neq0\)) на этом рисунке и на рис.~\cref{fig:insert_diag} ({\color{red}красные} кривые) были рассчитаны численно\cite{babichev2013black, dokuchaev2015physical,  dokuchaev2017gravitational, dokuchaev2019event_a, dokuchaev2019event_b, dokuchaevSeeInvisibleImage2019} с использованием соответствующих уравнений движения в метрике Керра–Ньюмена \ref{eq:EoM}.
%
%В то же время это изображение вводит в заблуждение и физически противоречиво: действительно, путешествие начинается в апогее \(r_a\) из точки \({\color{magenta}A}\), достигает перигея \(r_p\) и возвращается в апогей \(r_a\) в точке \(B\) за конечное собственное время, демонстрируя абсолютную геометрическую симметрию. Однако есть критический нюанс: этот апогей \(r_a\) в точке \(B\) находится не в родной вселенной, а в другой, весьма удаленной вселенной, что ясно видно на диаграмме Картера–Пенроуза на рис. \cref{fig:Pen_diag}. Этот нюанс снова разрушает симметрию моста Эйнштейна–Розена.
%
%На рис.~\cref{fig:insert_diag} показана диаграмма вложения для путешествия через внутренности черной дыры с использованием моста Эйнштейна–Розена. Диаграмма вложения очень полезна для интуитивного понимания особенностей загадочных черных дыр. На этой диаграмме мост Эйнштейна–Розена соединяет две асимптотически плоские вселенные, подобно туннелю кротовой норы \cite{morris1988wormholes}, но только с односторонним движением от начальной точки \(A\) к конечной точке \(B\). Геометрическая { симметрия} этой диаграммы вложения обманчива. На самом деле, она демонстрирует { асимметричную} пространственно-временную природу одностороннего моста Эйнштейна–Розена (с учетом потери конуса).

%Полностью симметричная картина периодического орбитального движения пробной планеты или космического корабля вокруг центральной сингулярности черной дыры Рейсснера–Нордстрёма внутри горизонта Коши \(r_-\) показана на рис. \cref{fig:2D_traj_b}. Электрический заряд черной дыры составляет \(e=0.99\), а пробная планета (или корабль) с зарядом \(\varepsilon=-1.5\) периодически обращается вокруг черной дыры с параметрами \(\gamma=1.7\), \(\lambda=0.1\), соответствующими максимальному радиусу \(r_{ max}=0.75\) и минимальному радиусу \(r_{ man}=0.09\). Асимметричный мост Рейсснера–Нордстрёма необходим только для проникновения в эту очень экзотическую область при \(0<r<r_-\), где существуют почти стабильные периодические орбиты для пробных частиц \cite{grunauGeodesicsElectricallyMagnetically2011, olivares2011motion, puglieseEquatorialCircularMotion2011}, очень похожие на периодические орбиты вне горизонта событий \(r_+\).



\section{Проблема инфляции массы}\label{sec:2.3}

\subsection{Постановка проблемы}\label{sub:2.3.1}

При рассмотрении движения частиц в метрике черной дыры мы исходили из предположения малого влияния частицы на геометрию пространства и полностью им пренебрегали, однако для более детального изучения процесса прохождения наблюдателя через горизонт также необходимо учесть изменение самой метрики черной дыры после поглощения вещества. Для этого необходимо решить уравнения Эйнштейна с учетом падающего вещества, для чего можно использовать метрику Вайдья (формула \cref{eq:Vaidya}) \cite{vaidya1951nonstatic, vaidya1951gravitational}.

Проблема инфляции массы возникает при анализе внутренней структуры заряженных (Райсснера-Нордстрёма) и вращающихся (Керра) черных дыр, где присутствует внутренний горизонт Коши. Физически этот эффект проявляется как экспоненциальный рост энергии возмущений при приближении к внутреннему горизонту, описываемый соотношением \(\delta m \sim m_0 e^{k_-v}\), где \(k_-\) — поверхностная гравитация внутреннего горизонта, v — световая координата (формула \cref{eq:v}). Теоретический анализ показывает, что любое сколь угодно малое возмущение метрики (например, падающая частица или квантовая флуктуация) приводит к катастрофическому увеличению локальной кривизны, превращая внутренний горизонт в сингулярность пространственноподобного типа. Этот эффект ставит под сомнение физическую реализуемость внутренней структуры классических решений с горизонтом Коши.

Строгий математический анализ проблемы инфляции массы был впервые проведен Пуассоном и Исраэлем в 1989 году \cite{poissonInnerhorizonInstabilityMass1989, poissonEschatologyBlackHole1989}. В работе показано, что неустойчивость обусловлена комбинацией нескольких факторов: голубым смещением бесконечного порядка для фотонов, пересекающих горизонт Коши;  неконсервативной геометрией внутренней области, допускающей накопление энергии; и положительной обратной связью между возмущениями и метрикой. Уравнения возмущений демонстрируют режим с обострением, при котором скалярная кривизна \(R\) растет как \(R ~ e^{2kv}\) \cite{poissonInternalStructureBlack1990}. Последующие исследования уточнили, что этот процесс происходит за конечное собственное время наблюдателя, падающего в черную дыру, что делает проблему ограничением классической ОТО.

\subsection{Возможное разрешение проблемы}\label{sub:2.3.2}

В оригинальных работах \cite{poissonInnerhorizonInstabilityMass1989, poissonEschatologyBlackHole1989} приводятся аналитические вычисления для случая падающей сферической бесконечно-тонкой оболочки с использованием метрики Вайдья (формула \cref{eq:Vaidya}), однако такой случай соответствует бесконечной плотности энергии на самой оболочке, что приводит к скачкообразному изменению метрики за ней. Возможно, что в этом случае нет учета плавного изменения расположения горизонта Коши, на котором и возникает проблема инфляции массы \cite{dokuchaevMassInflationBlack2014}. В связи с тем, что горизонт с необходимостью должен смещаться при плавном изменении массы черной дыры, он перестает быть светоподобной поверхностью и становится пространственно подобной. Процесс плавного изменения горизонта Коши в случае излучения, распределенного в пространстве, показан на диаграмме Пенроуза (рис. \cref{fig:Pen_diag_2}.

\begin{figure}[ht]
	\centerfloat{
	\hfill
	\includegraphics[angle=0,width=0.6\linewidth]{Pen_diag_2} 
	\hfill
	}
	\caption[Диаграмма Картера-Пенроуза для метрики с излучением]{Диаграмма Картера-Пенроуза для сферически симметричной черной дыры и падающим на нее излучением} \label{fig:Pen_diag_2}
\end{figure}

Изучение этого случая в рамках классической ОТО является возможностью для дальнейшего развития концепции моста ЭР.

\subsection{Альтернативные подходы}\label{sub:2.3.1}

Проблема инфляции массы остается нерешенной задачей на стыке релятивистской астрофизики и квантовой гравитации, демонстрируя пределы применимости классической ОТО в экстремальных условиях. Ее решение требует разработки самосогласованных моделей, учитывающих квантовые эффекты в сильных гравитационных полях, и может привести к пересмотру современных представлений о конечной стадии гравитационного коллапса. Перспективные направления исследований включают построение квантовых моделей внутренней структуры черных дыр; поиск наблюдательных проявлений в данных гравитационно-волновых детекторов и радиоинтерферометров; изучение связи с проблемой информации и голографическим принципом. Особый интерес представляет разработка новых теоретических подходов, объединяющих методы квантовой теории поля, петлевой квантовой гравитации и струнных теорий, которые могли бы обеспечить регуляризацию расходимостей на внутреннем горизонте. Решение этой проблемы не только углубит понимание структуры пространства-времени в экстремальных условиях, но и послужит важным тестом для альтернативных теорий гравитации, потенциально открывая путь к созданию последовательной теории квантовой гравитации.

%\section{Одиночное изображение}\label{sec:ch2/sec1}
%
%\begin{figure}[ht]
%    \centerfloat{
%        \includegraphics[scale=0.27]{latex}
%    }
%    \caption{TeX.}\label{fig:latex}
%\end{figure}

%Для выравнивания изображения по-центру используется команда \verb+\centerfloat+, которая является во
%многом улучшенной версией встроенной команды \verb+\centering+.
%
%\section{Длинное название параграфа, в котором мы узнаём как сделать две картинки с~общим номером и названием}\label{sec:ch2/sect2}
%
%А это две картинки под общим номером и названием:
%\begin{figure}[ht]
%    \begin{minipage}[b][][b]{0.49\linewidth}\centering
%        \includegraphics[width=0.5\linewidth]{knuth1} \\ а)
%    \end{minipage}
%    \hfill
%    \begin{minipage}[b][][b]{0.49\linewidth}\centering
%        \includegraphics[width=0.5\linewidth]{knuth2} \\ б)
%    \end{minipage}
%    \caption{Очень длинная подпись к изображению,
%        на котором представлены две фотографии Дональда Кнута}
%    \label{fig:knuth}
%\end{figure}
%
%Те~же~две картинки под~общим номером и~названием,
%но с автоматизированной нумерацией подрисунков:
%\begin{figure}[ht]
%    \centerfloat{
%        \hfill
%        \subcaptionbox[List-of-Figures entry]{Первый подрисунок\label{fig:knuth_2-1}}{%
%            \includegraphics[width=0.25\linewidth]{knuth1}}
%        \hfill
%        \subcaptionbox{\label{fig:knuth_2-2}}{%
%            \includegraphics[width=0.25\linewidth]{knuth2}}
%        \hfill
%        \subcaptionbox{Третий подрисунок, подпись к которому
%            не~помещается на~одной строке}{%
%            \includegraphics[width=0.3\linewidth]{example-image-c}}
%        \hfill
%    }
%    \legend{Подрисуночный текст, описывающий обозначения, например. Согласно
%        ГОСТ 2.105, пункт 4.3.1, располагается перед наименованием рисунка.}
%    \caption[Этот текст попадает в названия рисунков в списке рисунков]{Очень
%        длинная подпись к второму изображению, на~котором представлены две
%        фотографии Дональда Кнута}\label{fig:knuth_2}
%\end{figure}
%
%На рисунке~\cref{fig:knuth_2-1} показан Дональд Кнут без головного убора.
%На рисунке~\cref{fig:knuth_2}\subcaptionref*{fig:knuth_2-2}
%показан Дональд Кнут в головном уборе.
%
%\section{Векторная графика}\label{sec:ch2/vector}
%
%Возможно вставлять векторные картинки, рассчитываемые \LaTeX\ <<на~лету>>
%с~их~предварительной компиляцией. Надписи в таких рисунках будут выполнены
%тем же~шрифтом, который указан для документа в целом.
%На~рисунке~\cref{fig:tikz_example} на~странице~\pageref{fig:tikz_example}
%представлен пример схемы, рассчитываемой пакетом \verb|tikz| <<на~лету>>.
%Для ускорения компиляции, подобные рисунки могут быть <<кешированы>>, что
%определяется настройками в~\verb|common/setup.tex|.
%Причём имя предкомпилированного
%файла и~папка расположения таких файлов могут быть отдельно заданы,
%что удобно, если не~для подготовки диссертации,
%то~для подготовки научных публикаций.
%\begin{figure}[ht]
%    \centerfloat{
%        \ifdefmacro{\tikzsetnextfilename}{\tikzsetnextfilename{tikz_example_compiled}}{}% присваиваемое предкомпилированному pdf имя файла (не обязательно)
%        \input{Dissertation/images/tikz_scheme.tikz}
%
%    }
%    \legend{}
%    \caption[Пример \texttt{tikz} схемы]{Пример рисунка, рассчитываемого
%        \texttt{tikz}, который может быть предкомпилирован}\label{fig:tikz_example}
%\end{figure}
%
%Множество программ имеют либо встроенную возможность экспортировать векторную
%графику кодом \verb|tikz|, либо соответствующий пакет расширения.
%Например, в GeoGebra есть встроенный экспорт,
%для Inkscape есть пакет svg2tikz,
%для Python есть пакет tikzplotlib,
%для R есть пакет tikzdevice.
%
%\begin{figure}[htbp]
%    \centerfloat{
%        \ifdefmacro{\tikzsetnextfilename}{\tikzsetnextfilename{pic2}}{}%
%        \input{Dissertation/images/scheme.tikz}
%    }
%    \legend{%
%        \textbf{1} "--- кружок с загогулиной;
%        \textbf{2} "--- камертоны;
%        \textbf{3} "--- кресты;
%        \textbf{4} "--- волны;
%        \textbf{5} "--- прямоугольники;
%        \textbf{5} "--- пронзённый стрелой прямоугольник.%
%    }
%    \caption{Составная схема \textit{tikz}}\label{fig:scheme-tikz}
%\end{figure}
%
%На рисунке~\cref{fig:scheme-tikz} представлена составная схема \textit{tikz}.
%Каждый её элемент нарисован в отдельном файле в единичном масштабе.
%Расстановка элементов на~рисунке производится при помощи аргументов \texttt{xshift},
%\texttt{yshift}, \texttt{rotate} и~\texttt{scale} окружения \texttt{scope}.
%
%Пример использования библиотеки \textit{circuitikz} изображён на рисунке~\cref{fig:circuitikz}.
%
%\begin{figure}[htbp]
%    \centerfloat{
%        \input{Dissertation/images/circuit.tikz}
%    }
%    \caption{Схема \textit{circuitikz}}\label{fig:circuitikz}
%\end{figure}
%
%Красивые графики также можно добавлять при помощи пакета \textit{pgfplot}~(рисунок~\cref{fig:pgfplot}).
%Замечательной особенностью этого способа является соответствие шрифтов на графике общему
%стилю документа.
%
%\begin{figure}[htbp]
%    \centerfloat{
%        \input{Dissertation/images/plot_csv.tikz}
%    }
%    \caption{График \textit{pgfplot} на основе данных из \texttt{csv} файла}\label{fig:pgfplot}
%\end{figure}
%
%
%\section{Пример вёрстки списков}\label{sec:ch2/sec3}
%
%\noindent Нумерованный список:
%\begin{enumerate}
%    \item Первый пункт.
%    \item Второй пункт.
%    \item Третий пункт.
%\end{enumerate}
%
%\noindent Маркированный список:
%\begin{itemize}
%    \item Первый пункт.
%    \item Второй пункт.
%    \item Третий пункт.
%\end{itemize}
%
%\noindent Вложенные списки:
%\begin{itemize}
%    \item Имеется маркированный список.
%          \begin{enumerate}
%              \item В нём лежит нумерованный список,
%              \item в котором
%                    \begin{itemize}
%                        \item лежит ещё один маркированный список.
%                    \end{itemize}
%          \end{enumerate}
%\end{itemize}
%
%\noindent Нумерованные вложенные списки:
%\begin{enumerate}
%    \item Первый пункт.
%    \item Второй пункт.
%    \item Вообще, по ГОСТ 2.105 первый уровень нумерации
%          (при необходимости ссылки в тексте документа на одно из перечислений)
%          идёт буквами русского или латинского алфавитов,
%          а второй "--- цифрами со~скобками.
%          Здесь отходим от ГОСТ.
%          \begin{enumerate}
%              \item в нём лежит нумерованный список,
%              \item в котором
%                    \begin{enumerate}
%                        \item ещё один нумерованный список,
%                        \item третий уровень нумерации не нормирован ГОСТ 2.105;
%                        \item обращаем внимание на строчность букв,
%                        \item в этом списке
%                              \begin{itemize}
%                                  \item лежит ещё один маркированный список.
%                              \end{itemize}
%                    \end{enumerate}
%
%          \end{enumerate}
%
%    \item Четвёртый пункт.
%\end{enumerate}
%
%\section{Традиции русского набора}
%
%Много полезных советов приведено в материале
%<<\href{https://kostyrka.ru/main/ru/typesetting-and-typography-crash-course-by-kostyrka/}{Краткий курс благородного набора}>>
%(автор А.\:В.~Костырка).
%Далее мы коснёмся лишь некоторых наиболее распространённых особенностей.
%
%\subsection{Пробелы}
%
%В~русском наборе принято:
%\begin{itemize}
%    \item единицы измерения, знак процента отделять пробелами от~числа:
%          10~кВт, 15~\% (согласно ГОСТ 8.417, раздел 8);
%    \item \(\tg 20\text{\textdegree}\), но: 20~{\textdegree}C
%          (согласно ГОСТ 8.417, раздел 8);
%    \item знак номера, параграфа отделять от~числа: №~5, \S~8;
%    \item стандартные сокращения: т.\:е., и~т.\:д., и~т.\:п.;
%    \item неразрывные пробелы в~предложениях.
%\end{itemize}
%
%\subsection{Математические знаки и символы}
%
%Русская традиция начертания греческих букв и некоторых математических
%функций отличается от~западной. Это исправляется серией
%\verb|\renewcommand|.
%\begin{itemize}
%    %Все \original... команды заранее, ради этого примера, определены в Dissertation\userstyles.tex
%    \item[До:] \( \originalepsilon \originalge \originalphi\),
%          \(\originalphi \originalleq \originalepsilon\),
%          \(\originalkappa \in \originalemptyset\),
%          \(\originaltan\),
%          \(\originalcot\),
%          \(\originalcsc\).
%    \item[После:] \( \epsilon \ge \phi\),
%          \(\phi \leq \epsilon\),
%          \(\kappa \in \emptyset\),
%          \(\tan\),
%          \(\cot\),
%          \(\csc\).
%\end{itemize}
%
%Кроме того, принято набирать греческие буквы вертикальными, что
%решается подключением пакета \verb|upgreek| (см. закомментированный
%блок в~\verb|userpackages.tex|) и~аналогичным переопределением в
%преамбуле (см.~закомментированный блок в~\verb|userstyles.tex|). В
%этом шаблоне такие переопределения уже включены.
%
%Знаки математических операций принято переносить. Пример переноса
%в~формуле~\eqref{eq:equation3}.
%
%\subsection{Кавычки}
%В английском языке приняты одинарные и двойные кавычки в~виде ‘...’ и~“...”.
%В~России приняты французские («...») и~немецкие („...“) кавычки (они называются
%«ёлочки» и~«лапки», соответственно). ,,Лапки`` обычно используются внутри
%<<ёлочек>>, например, <<... наш гордый ,,Варяг``...>>.
%
%Французкие левые и правые кавычки набираются
%как лигатуры \verb|<<| и~\verb|>>|, а~немецкие левые
%и правые кавычки набираются как лигатуры \verb|,,| и~\verb|‘‘| (\verb|``|).
%
%Вместо лигатур или команд с~активным символом "\ можно использовать команды
%\verb|\glqq| и \verb|\grqq| для набора немецких кавычек и команды \verb|\flqq|
%и~\verb|\frqq| для набора французских кавычек. Они определены в пакете
%\verb|babel|.
%
%\subsection{Тире}
%%  babel+pdflatex по умолчанию, в polyglossia надо включать опцией (и перекомпилировать с удалением временных файлов)
%Команда \verb|"---| используется для печати тире в тексте. Оно может быть
%несколько короче английского длинного тире (подробности в~документации
%русификации babel). Кроме того, команда задаёт небольшую жёсткую отбивку
%от~слова, стоящего перед тире. При этом, само тире не~отрывается от~слова.
%После тире следует такая же отбивка от текста, как и~перед тире. При наборе
%текста между словом и командой, за которым она следует, должен стоять пробел.
%
%В составных словах, таких, как <<Закон Менделеева"--~Клапейрона>>, для печати
%тире надо использовать команду \verb|"--~|. Она ставит более короткое,
%по~сравнению с~английским, тире и позволяет делать переносы во втором слове.
%При~наборе текста команда \verb|"--~| не отделяется пробелом от слова,
%за~которым она следует (\verb|Менделеева"--~|). Следующее за командой слово
%может быть  отделено от~неё пробелом или перенесено на другую строку.
%
%Если прямая речь начинается с~абзаца, то перед началом её печатается тире
%командой \verb|"--*|. Она печатает русское тире и жёсткую отбивку нужной
%величины перед текстом.
%
%\subsection{Дефисы и переносы слов}
%%  babel+pdflatex по умолчанию, в polyglossia надо включать опцией (и перекомпилировать с удалением временных файлов)
%Для печати дефиса в~составных словах введены две команды. Команда~\verb|"~|
%печатает дефис и~запрещает делать переносы в~самих словах, а~команда \verb|"=|
%печатает дефис, оставляя \TeX ’у право делать переносы в~самих словах.
%
%В отличие от команды \verb|\-|, команда \verb|"-| задаёт место в~слове, где
%можно делать перенос, не~запрещая переносы и~в~других местах слова.
%
%Команда \verb|""| задаёт место в~слове, где можно делать перенос, причём дефис
%при~переносе в~этом месте не~ставится.
%
%Команда \verb|",| вставляет небольшой пробел после инициалов с~правом переноса
%в~фамилии.
%
%\section{Текст из панграмм и формул}
%
%Любя, съешь щипцы, "--- вздохнёт мэр, "--- кайф жгуч. Шеф взъярён тчк щипцы
%с~эхом гудбай Жюль. Эй, жлоб! Где туз? Прячь юных съёмщиц в~шкаф. Экс-граф?
%Плюш изъят. Бьём чуждый цен хвощ! Эх, чужак! Общий съём цен шляп (юфть) "---
%вдрызг! Любя, съешь щипцы, "--- вздохнёт мэр, "--- кайф жгуч. Шеф взъярён тчк
%щипцы с~эхом гудбай Жюль. Эй, жлоб! Где туз? Прячь юных съёмщиц в~шкаф.
%Экс-граф? Плюш изъят. Бьём чуждый цен хвощ! Эх, чужак! Общий съём цен шляп
%(юфть) "--- вдрызг! Любя, съешь щипцы, "--- вздохнёт мэр, "--- кайф жгуч. Шеф
%взъярён тчк щипцы с~эхом гудбай Жюль. Эй, жлоб! Где туз? Прячь юных съёмщиц
%в~шкаф. Экс-граф? Плюш изъят. Бьём чуждый цен хвощ! Эх, чужак! Общий съём цен
%шляп (юфть) "--- вдрызг! Любя, съешь щипцы, "--- вздохнёт мэр, "--- кайф жгуч.
%Шеф взъярён тчк щипцы с~эхом гудбай Жюль. Эй, жлоб! Где туз? Прячь юных съёмщиц
%в~шкаф. Экс-граф? Плюш изъят. Бьём чуждый цен хвощ! Эх, чужак! Общий съём цен
%шляп (юфть) "--- вдрызг! Любя, съешь щипцы, "--- вздохнёт мэр, "--- кайф жгуч.
%Шеф взъярён тчк щипцы с~эхом гудбай Жюль. Эй, жлоб! Где туз? Прячь юных съёмщиц
%в~шкаф. Экс-граф? Плюш изъят. Бьём чуждый цен хвощ! Эх, чужак! Общий съём цен
%шляп (юфть) "--- вдрызг! Любя, съешь щипцы, "--- вздохнёт мэр, "--- кайф жгуч.
%Шеф взъярён тчк щипцы с~эхом гудбай Жюль. Эй, жлоб! Где туз? Прячь юных съёмщиц
%в~шкаф. Экс-граф? Плюш изъят. Бьём чуждый цен хвощ! Эх, чужак! Общий съём цен
%шляп (юфть) "--- вдрызг! Любя, съешь щипцы, "--- вздохнёт мэр, "--- кайф жгуч.
%Шеф взъярён тчк щипцы с~эхом гудбай Жюль. Эй, жлоб! Где туз? Прячь юных съёмщиц
%в~шкаф. Экс-граф? Плюш изъят. Бьём чуждый цен хвощ! Эх, чужак! Общий съём цен
%шляп (юфть) "--- вдрызг! Любя, съешь щипцы, "--- вздохнёт мэр, "--- кайф жгуч.
%Шеф взъярён тчк щипцы с~эхом гудбай Жюль. Эй, жлоб! Где туз? Прячь юных съёмщиц
%в~шкаф. Экс-граф? Плюш изъят. Бьём чуждый цен хвощ! Эх, чужак! Общий съём цен
%шляп (юфть) "--- вдрызг! Любя, съешь щипцы, "--- вздохнёт мэр, "--- кайф жгуч.
%Шеф взъярён тчк щипцы с~эхом гудбай Жюль. Эй, жлоб! Где туз? Прячь юных съёмщиц
%в~шкаф. Экс-граф? Плюш изъят. Бьём чуждый цен хвощ! Эх, чужак! Общий съём цен
%шляп (юфть) "--- вдрызг! Любя, съешь щипцы, "--- вздохнёт мэр, "--- кайф жгуч.
%Шеф взъярён тчк щипцы с~эхом гудбай Жюль. Эй, жлоб! Где туз? Прячь юных съёмщиц
%в~шкаф. Экс-граф? Плюш изъят. Бьём чуждый цен хвощ! Эх, чужак! Общий съём цен
%шляп (юфть) "--- вдрызг! Любя, съешь щипцы, "--- вздохнёт мэр, "--- кайф жгуч.
%Шеф взъярён тчк щипцы с~эхом гудбай Жюль. Эй, жлоб! Где туз? Прячь юных съёмщиц
%в~шкаф. Экс-граф? Плюш изъят. Бьём чуждый цен хвощ! Эх, чужак! Общий съём цен
%шляп (юфть) "--- вдрызг!Любя, съешь щипцы, "--- вздохнёт мэр, "--- кайф жгуч.
%Шеф взъярён тчк щипцы с~эхом гудбай Жюль. Эй, жлоб! Где туз? Прячь юных съёмщиц
%в~шкаф. Экс-граф? Плюш изъят. Бьём чуждый цен хвощ! Эх, чужак! Общий съём цен
%
%Ку кхоро адолэжкэнс волуптариа хаж, вим граэко ыкчпэтында ты. Граэкы жэмпэр
%льюкяльиюч квуй ку, аэквюы продыжщэт хаж нэ. Вим ку магна пырикульа, но квюандо
%пожйдонёюм про. Квуй ат рыквюы ёнэрмйщ. Выро аккузата вим нэ.
%\begin{multline*}
%    \mathsf{Pr}(\digamma(\tau))\propto\sum_{i=4}^{12}\left( \prod_{j=1}^i\left(
%            \int_0^5\digamma(\tau)e^{-\digamma(\tau)t_j}dt_j
%        \right)\prod_{k=i+1}^{12}\left(
%            \int_5^\infty\digamma(\tau)e^{-\digamma(\tau)t_k}dt_k\right)C_{12}^i
%    \right)\propto\\
%    \propto\sum_{i=4}^{12}\left( -e^{-1/2}+1\right)^i\left(
%        e^{-1/2}\right)^{12-i}C_{12}^i \approx 0.7605,\quad
%    \forall\tau\neq\overline{\tau}
%\end{multline*}
%Квуй ыёюз омниюм йн. Экз алёквюам кончюлату квуй, ты альяквюам ёнвидюнт пэр.
%Зыд нэ коммодо пробатуж. Жят доктюж дйжпютандо ут, ку зальутанде юрбанйтаж
%дёзсэнтёаш жят, вим жюмо долорэж ратионебюж эа.
%
%Ад ентэгры корпора жплэндидэ хаж. Эжт ат факэтэ дычэрунт пэржыкюти. Нэ нам
%доминг пэрчёус. Ку квюо ёужто эррэм зючкёпит. Про хабэо альбюкиюс нэ.
%\[
%    \begin{pmatrix}
%        a_{11} & a_{12} & a_{13} \\
%        a_{21} & a_{22} & a_{23}
%    \end{pmatrix}
%\]
%
%\[
%    \begin{vmatrix}
%        a_{11} & a_{12} & a_{13} \\
%        a_{21} & a_{22} & a_{23}
%    \end{vmatrix}
%\]
%
%\[
%    \begin{bmatrix}
%        a_{11} & a_{12} & a_{13} \\
%        a_{21} & a_{22} & a_{23}
%    \end{bmatrix}
%\]
%Про эа граэки квюаыквуэ дйжпютандо. Ыт вэл тебиквюэ дэфянятйоныс, нам жолюм
%квюандо мандамюч эа. Эож пауло лаудым инкедыринт нэ, пэрпэтюа форынчйбюж пэр
%эю. Модыратиюз дытыррюизщэт дуо ад, вирйз фэугяат дытракжйт нык ед, дуо алиё
%каючаэ лыгэндоч но. Эа мольлиз юрбанйтаж зигнёфэрумквюы эжт.
%
%Про мандамюч кончэтытюр ед. Трётанё прёнкипыз зигнёфэрумквюы вяш ан. Ат хёз
%эквюедым щуавятатэ. Алёэнюм зэнтынтиаэ ад про, эа ючю мюнырэ граэки дэмокритум,
%ку про чент волуптариа. Ыльит дыкоры аляквюид еюж ыт. Ку рыбюм мюндй ютенам
%дуо.
%\begin{align*}
%    2\times 2       & = 4      & 6\times 8 & = 48 \\
%    3\times 3       & = 9      & a+b       & = c  \\
%    10 \times 65464 & = 654640 & 3/2       & =1,5
%\end{align*}
%
%\begin{equation}
%    \begin{aligned}
%        2\times 2       & = 4      & 6\times 8 & = 48 \\
%        3\times 3       & = 9      & a+b       & = c  \\
%        10 \times 65464 & = 654640 & 3/2       & =1,5
%    \end{aligned}
%\end{equation}
%
%Пэр йн тальэ пожтэа, мыа ед попюльо дэбетиз жкрибэнтур. Йн квуй аппэтырэ
%мэнандря, зыд аляквюид хабымуч корпора йн. Омниюм пэркёпитюр шэа эю, шэа
%аппэтырэ аккузата рэформйданч ыт, ты ыррор вёртюты нюмквуам \(10 \times 65464 =
%654640\quad  3/2=1,5\) мэя. Ипзум эуежмод \(a+b = c\) мальюизчыт ад дуо. Ад
%фэюгаят пытынтёюм адвыржаряюм вяш. Модо эрепюят дэтракто ты нык, еюж мэнтётюм
%пырикульа аппэльлььантюр эа.
%
%Мэль ты дэлььынётё такематыш. Зэнтынтиаэ конклььюжионэмквуэ ан мэя. Вёжи лебыр
%квюаыквуэ квуй нэ, дуо зймюл дэлььиката ку. Ыам ку алиё путынт.
%
%%Большая фигурная скобка только справа
%\[\left. %ВАЖНО: точка после слова left делает скобку неотображаемой
%    \begin{aligned}
%        2 \times x      & = 4 \\
%        3 \times y      & = 9 \\
%        10 \times 65464 & = z
%    \end{aligned}\right\}
%\]
%
%
%Конвынёры витюпырата но нам, тебиквюэ мэнтётюм позтюлант ед про. Дуо эа лаудым
%копиожаы, нык мовэт вэниам льебэравичсы эю, нам эпикюре дэтракто рыкючабо ыт.
%Вэрйтюж аккюжамюз ты шэа, дэбетиз форынчйбюж жкряпшэрит ыт прё. Ан еюж тымпор
%рыфэррэнтур, ючю дольор котёдиэквюэ йн. Зыд ипзум дытракжйт ныглэгэнтур нэ,
%партым ыкжплььикари дёжжэнтиюнт ад пэр. Мэль ты кытэрож молыжтйаы, нам но ыррор
%жкрипта аппарэат.
%
%\[ \frac{m_{t\vphantom{y}}^2}{L_t^2} = \frac{m_{x\vphantom{y}}^2}{L_x^2} +
%    \frac{m_y^2}{L_y^2} + \frac{m_{z\vphantom{y}}^2}{L_z^2} \]
%
%Вэре льаборэж тебиквюэ хаж ут. Ан пауло торквюатоз хаж, нэ пробо фэугяат
%такематыш шэа. Мэльёуз пэртинакёа юлламкорпэр прё ад, но мыа рыквюы конкыптам.
%Хёз квюот пэртинакёа эи, ельлюд трактатоз пэр ад. Зыд ед анёмал льаборэж
%номинави, жят ад конгуы льабятюр. Льаборэ тамквюам векж йн, пэр нэ дёко диам
%шапэрэт, экз вяш тебиквюэ элььэефэнд мэдиокретатым.
%
%Нэ про натюм фюйзчыт квюальизквюэ, аэквюы жкаывола мэль ку. Ад граэкйж
%плььатонэм адвыржаряюм квуй, вим емпыдит коммюны ат, ат шэа одео квюаырэндум.
%Вёртюты ажжынтиор эффикеэнди эож нэ, доминг лаборамюз эи ыам. Чэнзэрет
%мныжаркхюм экз эож, ыльит тамквюам факильизиж нык эи. Квуй ан элыктрам
%тинкидюнт ентырпрытаряш. Йн янвыняры трактатоз зэнтынтиаэ зыд. Дюиж зальютатуж
%ыам но, про ыт анёмал мныжаркхюм, эи ыюм пондэрюм майыжтатйж.
%
%\FloatBarrier
